% Table: Case study — teacher vs student reasoning traces.
\begin{table}[t]
\centering
\caption{Representative case study comparing teacher (Qwen3-32B + TopoPRM) and
  student (Qwen3-8B + RKL) reasoning traces on a middle-school critique problem.
  The student produces more concise reasoning while preserving key logical steps.
  Full traces are provided in Appendix~\ref{app:case_study}.}
\label{tab:case_study}
\vspace{2pt}
\resizebox{\textwidth}{!}{%
\begin{tabular}{@{}p{2.2cm} p{5.5cm} p{5.5cm}@{}}
\toprule
& \textbf{Teacher (32B)} & \textbf{Student (8B)} \\
\midrule
Problem & \multicolumn{2}{c}{\textit{(Example problem to be filled after experiment)}} \\
\midrule
\# Steps & TBD & TBD \\
\# Tokens & TBD & TBD \\
Score-Acc & TBD & TBD \\
\midrule
Key reasoning & TBD & TBD \\
\bottomrule
\end{tabular}%
}
\end{table}
